\subsection{SFT 数据系统：少量样本为何能快速塑形}

slides 30--36 先给出 SFT 的基本机制：在少量理想回答上继续做 language modeling。人工示范曾是 ChatGPT 式指令模型的关键，但成本和速度限制促使 Alpaca 使用强模型生成指令数据。若有 verifier，还可以让临时模型生成多个候选，再通过测试、规则或偏好筛选正确样本。Kimi K2 展示了更复杂的 Agent 数据系统：模拟用户和工具、生成轨迹，再用 rubric 和 grounded code 信号过滤。

\lecturefigure{slides-images/slide-030.png}{官方 slide 30：alignment 任务与少量目标行为数据}
\lecturefigure{slides-images/slide-031.png}{官方 slide 31：SFT 对理想答案继续做 next-token prediction}
\lecturefigure{slides-images/slide-032.png}{官方 slide 32：Alpaca 用 LLM 扩展指令数据}
\lecturefigure{slides-images/slide-033.png}{官方 slide 33：DeepSeek-R1 式 verifier rejection sampling}
\lecturefigure{slides-images/slide-034.png}{官方 slide 34：SFT 可学习格式、工具使用与早期推理}
\lecturefigure{slides-images/slide-035.png}{官方 slide 35：Kimi K2 的模拟用户、工具与 rubric 数据流水线}
\lecturefigure{slides-images/slide-036.png}{官方 slide 36：SFT 数据量、LIMA 与 DeepSeek-R1 数量级}

\begin{knowledgebox}{读图：synthetic data 的关键是外部判据}
slide 32 的 teacher model 路径容易被误读成“让模型复制更强模型”。slide 33 增加 verifier 后，生成器只负责扩大候选空间，正确性由测试或其他外部规则决定；这降低了对 teacher 全面更强的要求。slide 35 又把任务推广到工具环境：模拟器必须产生可执行反馈，rubric 必须检查轨迹是否真正完成目标。数据系统的核心因此是 generator、environment、verifier 三者分工，而不是单一模型自举。
\end{knowledgebox}

\begin{warningbox}{SFT 学得快，不代表获得了新知识}
slide 36 的 LIMA 结果支持“少量高质量样本足以塑造格式与 instruction following”，因为大量知识已经在预训练中存在。它不支持用一万条样本从零教会基础模型一个完全陌生领域。对 Agent tool use，模型还可能只学会 JSON 外形，却没有学会何时调用、怎样处理异常和如何验证副作用。评估必须把格式正确与任务成功分开。
\end{warningbox}

\textbf{实践经验：}讲者在 01:02--01:08 特别推荐阅读 Kimi K2 的数据构造方法。这里的教学重点不是复刻某个私有 pipeline，而是学会把复杂能力拆成可生成、可执行、可筛选的轨迹。下一组 slides 转向 RL，正是因为行为克隆始终受示范者能力限制，并可能在模型不知道答案时教出“自信地编造”。

